\documentclass[11pt]{article}
\usepackage[T1]{fontenc}
\usepackage[utf8]{inputenc}
\usepackage{lmodern,amsmath,amssymb,amsthm}
\usepackage[margin=1in]{geometry}
\usepackage{newunicodechar}
\usepackage{hyperref}
\hypersetup{colorlinks=true,linkcolor=blue,citecolor=blue,urlcolor=blue,
  pdftitle={Symmetry of hook+column coefficients in iterated plethysms},pdfauthor={P. Graham}}
\newtheorem{theorem}{Theorem}
\newtheorem{corollary}[theorem]{Corollary}
\newtheorem{lemma}[theorem]{Lemma}
\newcommand{\Q}{\mathbb{Q}}
\newcommand{\Z}{\mathbb{Z}}
\newunicodechar{α}{\ensuremath{\alpha}}
\newunicodechar{β}{\ensuremath{\beta}}
\newunicodechar{γ}{\ensuremath{\gamma}}
\newunicodechar{ℕ}{\ensuremath{\mathbb{N}}}
\newunicodechar{ℚ}{\ensuremath{\mathbb{Q}}}
\newunicodechar{ℤ}{\ensuremath{\mathbb{Z}}}
\newunicodechar{∀}{\ensuremath{\forall}}
\newunicodechar{→}{\ensuremath{\to}}
\newunicodechar{₀}{\ensuremath{_0}}
\newunicodechar{≤}{\ensuremath{\leq}}
\newunicodechar{∧}{\ensuremath{\wedge}}
\newunicodechar{∨}{\ensuremath{\vee}}
\title{Symmetry of hook+column coefficients in iterated plethysms}
\author{P. Graham\\ \small\url{https://orcid.org/0009-0007-8493-1505}}
\date{November 1, 2026}
\begin{document}
\maketitle
\begin{abstract}
A hook+column partition $(\alpha,2^\beta,1^\gamma)$ has a first part $\alpha$, followed by $\beta$ parts equal to two and $\gamma$ parts equal to one. Let $g$ be a homogeneous rational symmetric function of degree $d\geq1$, let $a\geq2$, and put $f=g[h_a]$. We prove that, for each fixed $\gamma$, the Schur coefficients of $f$ at these partitions, listed by $\beta$ from $0$ to $d-1-\gamma$, read the same forwards and backwards, zero entries included, and that they vanish for larger $\beta$. Consequently, every iterated plethysm of complete homogeneous functions has symmetric hook+column sequences, answering Question 1 of Gutiérrez and Rosas. They gave explicit formulas for two three-factor families and proved that the second symmetric and exterior plethysms $s_{(2)}[f]$ and $s_{(1,1)}[f]$ of a Schur-positive $f$ preserve flip-symmetry, a stronger form of this symmetry; Langley and Remmel gave the underlying two-factor formula. We extend the symmetry to every number of factors and every positive index, and retain the zero entries. The argument uses Gutiérrez and Rosas's specialization at the signed alphabet $1-x-y$. After division by $(1-x)(1-y)$, replacing $x$ and $y$ by their reciprocals gives the coefficient identity. A Lean formalization proves the result from power-sum, Newton, and Jacobi--Trudi definitions, including the existence and uniqueness of Schur expansions.

This work has not yet been reviewed by experts. The proofs, the Lean formalization and the text were produced by AI models (from OpenAI and Anthropic), working under the author's direction.
\end{abstract}
\medskip
\noindent\textit{2020 Mathematics Subject Classification.} 05E05.\\
\textit{Keywords.} Plethysm, Schur functions, hook+column partitions.

\section{Introduction}
The identity $h_3[h_2]=s_{(6)}+s_{(4,2)}+s_{(2,2,2)}$ gives the hook+column coefficients $(b(0,0),b(1,0),b(2,0))=(1,1,1)$ in the notation defined below. We prove that the corresponding sequences are symmetric for arbitrary iterated plethysms of complete homogeneous functions.

Write $\Lambda_{\Q}=\Q[p_1,p_2,\ldots]$ for the ring of symmetric functions over $\Q$, generated by the power sums $p_j$, with $\deg p_j=j$. Let $h_n$ and $s_\lambda$ denote the complete homogeneous and Schur functions. The Schur functions form a basis of $\Lambda_{\Q}$, and $[s_\lambda]f$ denotes the coefficient of $s_\lambda$ when $f$ is expanded in this basis. The plethysm $f[g]$ is obtained from $f$ by replacing each $p_j$ with $p_j[g]$, and $p_j[g]$ is obtained from $g$ by replacing each $p_r$ with $p_{jr}$. As in~\cite[Definition 2.5]{GR}, $f\circ g=f[g]$, so $h_{n_1}\circ\cdots\circ h_{n_k}=h_{n_1}[h_{n_2}[\cdots[h_{n_k}]\cdots]]$, with $h_{n_k}$ innermost. A hook+column partition $(\alpha,2^\beta,1^\gamma)$ has a first part $\alpha$, followed by $\beta$ parts equal to two and $\gamma$ parts equal to one. As in Warning 2.2 of Gutiérrez and Rosas~\cite{GR}, the first part is counted separately even when it equals one or two; for example, $(2,2,2)$ has $(\alpha,\beta,\gamma)=(2,2,0)$. For a homogeneous function $f$ of degree $N$, define
\[
 b_f(\beta,\gamma)=[s_{(N-2\beta-\gamma,2^\beta,1^\gamma)}]f
 \qquad(\beta\in\Z,\ \gamma\geq0).
\]
Here the coefficient is zero unless $\beta\geq0$, $\alpha=N-2\beta-\gamma\geq1$, and either $\beta=0$ or $\alpha\geq2$. Thus $\alpha=1$ represents a column and occurs only when $\beta=0$.

\begin{theorem}\label{thm:main}
Let $g\in\Lambda_{\Q}$ be homogeneous of degree $d\geq1$, and let $a\geq2$ be an integer. Put $f=g[h_a]$, of degree $N=ad$. Then, for every $\beta\in\Z$ and $\gamma\geq0$,
\begin{equation}\label{eq:reflection}
 b_f(\beta,\gamma)=b_f(d-1-\gamma-\beta,\gamma).
\end{equation}
Moreover, $b_f(\beta,\gamma)=0$ if $\beta+\gamma\geq d$, and $[s_{(1^N)}]f=0$. No positivity condition on $g$ is required.
\end{theorem}

\begin{corollary}\label{cor:question}
Let $n_1,\ldots,n_k$ be a nonempty list of positive integers, put $f=h_{n_1}\circ\cdots\circ h_{n_k}$, and fix $\gamma\geq0$. The nonzero coefficients $b_f(\beta,\gamma)$, ordered by $\beta$, form a palindrome. More precisely, delete the indices equal to one. If none remain, put $d=1$; otherwise let $d$ be the product of the remaining indices except the last, so that $d=1$ when only one remains. Then $\beta\mapsto b_f(\beta,\gamma)$, for all $\beta\in\Z$ and with the zero convention above, satisfies \eqref{eq:reflection} and vanishes for $\beta+\gamma\geq d$.
\end{corollary}

This answers Question 1 of Gutiérrez and Rosas~\cite[Section 4]{GR}, posed in their January 2022 preprint and previously as Conjecture 1 in Gutiérrez's June 2021 thesis~\cite{thesis}. In their Definition 2.8, a sequence is symmetric if its nonzero entries form a symmetric sequence. Langley and Remmel~\cite[Theorem 4.8]{LR} give
\[
 (h_b[h_a])\big|_{\mathrm{hook+column}}
   =\sum_{j=0}^{b-1}s_{(ab-2j,2^j)}\quad(a,b\geq2),
\]
also stated as \cite[Theorem 2.14]{GR}. Theorems 2.15 and 2.16 of~\cite{GR} give explicit formulas for $h_2[h_b[h_a]]$ and $h_c[h_2[h_a]]$. Their Theorem 2.19 states that $s_{(2)}[f]$ and $s_{(1,1)}[f]$ preserve flip-symmetry, the stronger symmetry of their Definition 2.12, for Schur-positive $f$, that is, for $f$ with nonnegative Schur coefficients; the thesis contains the $s_{(2)}$ symmetry-preservation statement in Theorem 4.28. These are the prior partial results extended here.

Two further results of Gutiérrez and Rosas bear on the proof. The vanishing criterion behind Lemma~\ref{lem:schur} also appears in~\cite[Lemma 3.1]{positivity}: for alphabets $X_r$ and $Y_c$ of $r$ and $c$ variables, $s_\lambda[X_r-Y_c]$ is nonzero exactly when $\lambda$ is a $(c,r)$-hook, that is, when its $(r+1)$st part is at most $c$. Their Theorem~5.1 states that a nonzero coefficient of $s_\lambda$ in $s_\mu[s_\nu]$ forces $\nu\subseteq\lambda$; taking $\nu=(a)$ with $a\geq2$ and expanding $g$ in Schur functions, it already gives $[s_{(1^N)}]f=0$ in Theorem~\ref{thm:main}.

\paragraph{Proof overview.}
We use the signed-alphabet specialization and coefficient extraction of Gutiérrez and Rosas~\cite[Lemma 3.8]{GR}. At $1-x-y$, every $h_a$ with $a\geq2$ becomes $(1-x)(1-y)$. Weighted homogeneity then shows that replacing $x$ and $y$ by their reciprocals multiplies the specialized plethysm by $(xy)^{-d}$, where $d$ is the degree of $g$. Dividing by this common factor and multiplying by $x-y$ separates the hook+column coefficients into monomials above and below the diagonal. Reciprocal substitution exchanges those monomials and gives \eqref{eq:reflection}. The additional step beyond the cited specialization is this use of reciprocity for arbitrary $g$.

\section{Specialization and reflection}
Let $\phi:\Lambda_{\Q}\to\Q[x,y]$ be the algebra map $p_j\mapsto1-x^j-y^j$. Put $Q=(1-x)(1-y)$ and $u_\gamma(x,y)=\sum_{i=0}^\gamma x^i y^{\gamma-i}$.

\begin{lemma}\label{lem:schur}
One has $\phi(h_0)=1$, $\phi(h_1)=1-x-y$, and $\phi(h_n)=Q$ for $n\geq2$. For $\alpha\geq2$,
\begin{equation}\label{eq:schur}
 \phi(s_{(\alpha,2^\beta,1^\gamma)})
   =(-1)^\gamma Q(xy)^\beta u_\gamma(x,y).
\end{equation}
If $\lambda_2\geq3$, then $\phi(s_\lambda)=0$. Finally,
$\phi(s_{(1^N)})\big|_{x=1}=(-y)^N$.
\end{lemma}
\begin{proof}
The generating series for complete functions becomes
\[
 \sum_{n\geq0}\phi(h_n)t^n=\frac{(1-xt)(1-yt)}{1-t},
\]
which gives the first assertions. Apply these values to the Jacobi--Trudi matrix. When $\lambda_2\geq3$, its first two rows consist entirely of $Q$, so the determinant vanishes. For $(\alpha,2^\beta,1^\gamma)$, subtract each column's predecessor, proceeding from right to left. The first row becomes $(Q,0,\ldots,0)$. Expanding in this row leaves a block lower-triangular matrix: a lower-triangular block of size $\beta$ with diagonal $xy$ and zeros to its right, followed by a tridiagonal block of size $\gamma$ with diagonal $-(x+y)$, upper diagonal $xy$, and lower diagonal $1$. Its determinant is therefore $(xy)^\beta D_\gamma$, where the tridiagonal determinants satisfy
\[
 D_0=1,\quad D_1=-(x+y),\quad
 D_r=-(x+y)D_{r-1}-xyD_{r-2};
\]
thus $D_r=(-1)^r u_r$, proving \eqref{eq:schur}. At $x=1$, the column determinant has diagonal $-y$ and all entries above the diagonal zero, giving $(-y)^N$. This also proves directly the specialization in~\cite[Lemma 3.8]{GR}.
\end{proof}

\begin{lemma}\label{lem:recip}
Under the hypotheses of Theorem~\ref{thm:main}, let $P=\phi(g[h_a])$. Then $Q$ divides $P$, and, as a Laurent polynomial,
\[
 P(x^{-1},y^{-1})=(xy)^{-d}P(x,y).
\]
\end{lemma}
\begin{proof}
Plethysm substitutes $p_j[h_a]$ for each $p_j$ in $g$. Since $p_j[h_a]=h_a[p_j]$ is obtained from $h_a$ by the Adams operation $p_k\mapsto p_{jk}$, and $\phi(p_{jk})=1-(x^j)^k-(y^j)^k$, Lemma~\ref{lem:schur} gives $\phi(p_j[h_a])=\phi(h_a)(x^j,y^j)=(1-x^j)(1-y^j)$. Hence
\[
 P=g\big|_{p_j=(1-x^j)(1-y^j)}.
\]
Every monomial of $g$ has positive weighted degree and therefore contains a power-sum factor. Each substituted factor is divisible by $Q$, so $Q\mid P$. Under reciprocal substitution that factor becomes $(xy)^{-j}(1-x^j)(1-y^j)$. The sum of the indices, counted with multiplicity, is $d$ in each monomial of $g$, proving the identity.
\end{proof}

\begin{proof}[Proof of Theorem~\ref{thm:main}]
Expand $f$ in the Schur basis. By Lemma~\ref{lem:schur}, only hook+columns and the column survive specialization. Evaluating at $x=1$ kills $P$ by Lemma~\ref{lem:recip} and kills all hook+columns with first part at least two. Hence the column coefficient times $(-y)^N$ is zero. Since $\Q[y]$ is a domain, the column coefficient is zero.

Since $b_f(0,N-1)=[s_{(1^N)}]f=0$, every term with nonzero coefficient has $\alpha\geq2$, so~\eqref{eq:schur} applies to each term below. Write $G=P/Q$ and $H=(x-y)G$. The remaining terms give the finite sum
\begin{equation}\label{eq:H}
 H=\sum_{\beta,\gamma\geq0}(-1)^\gamma b_f(\beta,\gamma)
 \left(x^{\beta+\gamma+1}y^\beta-x^\beta y^{\beta+\gamma+1}\right).
\end{equation}
The monomials above the diagonal have distinct exponent pairs: if $i>j\geq0$, the coefficient of $x^iy^j$ is $(-1)^{i-j-1}b_f(j,i-j-1)$. Below the diagonal the coefficients have the opposite sign.

Since $Q(x^{-1},y^{-1})=(xy)^{-1}Q$, Lemma~\ref{lem:recip} yields
\begin{equation}\label{eq:Hrecip}
 H(x^{-1},y^{-1})=-(xy)^{-d}H(x,y).
\end{equation}
Extend the coefficients of $H$ by zero to all integer exponent pairs. Equation~\eqref{eq:Hrecip} says that the coefficients at $(i,j)$ and $(d-i,d-j)$ are negatives. In particular, a nonzero term has both exponents between zero and $d$.

Fix $\beta\geq0$ and $\gamma\geq0$. Its upper monomial in~\eqref{eq:H} has exponents $(\beta+\gamma+1,\beta)$. If $\beta+\gamma\geq d$, its first exponent exceeds $d$ and its coefficient is zero. Otherwise $\beta+\gamma\leq d-1$, so its reciprocal partner $(d-\beta-\gamma-1,\,d-\beta)$ has nonnegative exponents, lies below the diagonal and corresponds to
\[
 \beta'=d-1-\gamma-\beta\geq0.
\]
The sign in~\eqref{eq:Hrecip} cancels the minus sign below the diagonal in~\eqref{eq:H}, giving $b_f(\beta,\gamma)=b_f(\beta',\gamma)$. If $\beta\geq0$ and $\beta+\gamma\geq d$, both sides of~\eqref{eq:reflection} vanish: the left by the vanishing just shown, the right by definition, since $\beta'<0$. If instead $\beta<0$, the left side is zero by definition, while $\beta'+\gamma=d-1-\beta\geq d$, so the right side vanishes by the result just proved. This establishes all assertions.
\end{proof}

\begin{proof}[Proof of Corollary~\ref{cor:question}]
The function $h_1=p_1$ is the plethystic identity, and plethysm is associative. Delete factors equal to $h_1$. If any factors remain, take their last index as $a\geq2$ and their outer composite as $g$, using $p_1$ if there are no outer factors. Its degree is the stated product $d$. Theorem~\ref{thm:main} applies. For $\gamma<d$, the interval $0\leq\beta\leq d-1-\gamma$ is preserved by reflection, including all its zeros; for $\gamma\geq d$ every entry vanishes. Removing zeros therefore gives a palindrome. If all factors were $h_1$, the function is $p_1$; the sole nonzero hook+column coefficient is $b(0,0)=1$, and the conclusions hold with $d=1$.
\end{proof}

\section{Formal verification}
The formal theorem \texttt{GR22.fullEndpoint : GR22.FullEndpoint} is Theorem~\ref{thm:main}. Its statement and the coefficient it uses are
{\small
\begin{verbatim}
def FullEndpoint : Prop :=
  ∀ (d a : ℕ) (g : SF) (c : List ℕ →₀ ℚ),
    1 ≤ d → 2 ≤ a →
    MvPolynomial.IsWeightedHomogeneous (fun i : ℕ => i + 1) g d →
    SchurExpansion (a * d) (plethysm g (complete a)) c →
    (∀ (β : ℤ) (γ : ℕ),
      hookCoeff (a * d) c β γ = hookCoeff (a * d) c ((d : ℤ) - 1 - (γ : ℤ) - β) γ) ∧
    (∀ (β : ℤ) (γ : ℕ), (d : ℤ) ≤ β + (γ : ℤ) → hookCoeff (a * d) c β γ = 0) ∧
    c (List.replicate (a * d) 1) = 0

def hookCoeff (N : ℕ) (c : List ℕ →₀ ℚ) (β : ℤ) (γ : ℕ) : ℚ :=
  let α : ℤ := (N : ℤ) - 2 * β - (γ : ℤ)
  if 0 ≤ β ∧ 1 ≤ α ∧ (β = 0 ∨ 2 ≤ α) then
    c (α.toNat :: (List.replicate β.toNat 2 ++ List.replicate γ 1))
  else 0
\end{verbatim}
}
\noindent Here \texttt{SF} is the polynomial ring $\Q[p_1,p_2,\ldots]$, with variable $i$ standing for $p_{i+1}$, so the weighted-homogeneity hypothesis says that $g$ is homogeneous of degree $d$; the hypotheses $1\leq d$ and $2\leq a$ are those of Theorem~\ref{thm:main}. The term \texttt{plethysm g (complete a)} is $f=g[h_a]$, and \texttt{SchurExpansion (a * d) f c} says that $f=\sum_\mu c_\mu s_\mu$ over partitions $\mu$ of $N=ad$. The theorem \texttt{GR22.schurExpansion\_unique} shows that such a $c$ is unique, so $c_\mu=[s_\mu]f$. Then \texttt{hookCoeff N c β γ} is $b_f(\beta,\gamma)$: its guard is the convention of the introduction, and its list is the partition $(\alpha,2^\beta,1^\gamma)$. The three conjuncts are the reflection~\eqref{eq:reflection}, the vanishing for $\beta+\gamma\geq d$, and $[s_{(1^N)}]f=0$, the column being \texttt{List.replicate (a * d) 1}. The definitions use Newton's recurrence for $h_n$, Jacobi--Trudi for $s_\lambda$, and Adams substitution for plethysm.

The separate theorems \texttt{GR22.expansionsExist} and \texttt{GR22.schurExpansion\_unique} establish existence and uniqueness of these expansions from the definitions. Theorems \texttt{GR22.question1}, \texttt{GR22.iterated\_has\_expansion}, and \texttt{GR22.question1Positional} give Corollary~\ref{cor:question}, its expansion, and its version retaining zeros. Thus neither the specialization lemma nor the existence of expansions is an unproved hypothesis. These theorems depend only on Lean's standard axioms \texttt{propext}, \texttt{Classical.choice}, and \texttt{Quot.sound}.

The default build also checks the two identities $h_2[h_2]=s_4+s_{22}$ and $h_3[h_2]=s_6+s_{42}+s_{222}$. The latter refutes the erroneous reflection $\beta\mapsto d-\gamma-\beta$: at $d=3$, it would equate $b(0,0)=1$ with $b(3,0)=0$. The accompanying standalone Lake project gives the pinned dependencies and standard cache-download and build commands in its README. Outside the formal check is only the identification of these definitions with the usual ring of symmetric functions and of \texttt{hookCoeff} with $b^{\beta}_{f,\gamma}$ of~\cite[Definition 2.8]{GR}; no empirical coefficient calculation is a premise of the theorem.

\section{Related work}
The published version of~\cite{GR} still poses Question~1 in Section~4, and Gutiérrez's February 2024 ITMAIA talk slides~\cite{itmaia} say in their final section that the symmetry ``seems to hold for many more families''. In the sources cited in this section we found no statement of~\eqref{eq:reflection} or an equivalent identity. The closest are the following.

Gutiérrez, Martínez, Szwej and Wildon study plethysms of $\mathrm{SL}_2$-representations. In~\cite{lifts}, Lemma~3.6 gives symmetries of the Schur coefficients of the $k$-fold map, which repeats each variable $k$ times, and Proposition~5.1 is an $\mathrm{SL}_2(\mathbb{F})$-isomorphism for a Weyl module labeled by a hook partition. In~\cite{filtration}, a related paper, Theorem~1.1 proves the filtration of $\mathrm{SL}_2(\mathbb{F})$-plethysms proposed in~\cite[Conjecture~6.1]{lifts}. Gutiérrez, Orellana, Saliola, Schilling and Zabrocki~\cite[Theorem~5.10]{geometric} prove a reciprocity for the generating function of the two-row Schur coefficients of $s_\mu[h_n]$; it replaces the outer partition $\mu$ by its conjugate, whereas~\eqref{eq:reflection} keeps $g$ fixed. Lee~\cite[Theorem~1.1]{lee} interprets $\langle s_\lambda,s_\mu[h_r]\rangle$ only for $\lambda_1\leq r+1$. Sarkar~\cite[Theorem~6.4 and Corollary~7.2]{sarkar} studies first-row stability for outer $h_n$. Ma, Liu and Huang~\cite[Theorem~1.1]{mlh} study the coefficient of $s_{(3,2^{3m-2},1)}$ in $h_{2m}[e_3]$, whose inner function is the elementary symmetric function $e_3$. We also compared Fischer, Gangl, Gutiérrez and Kumari~\cite{fggk}, on Hall--Littlewood specializations of types $\mathrm{B}$, $\mathrm{C}$ and $\mathrm{BC}$; Gutiérrez and Krattenthaler~\cite{gk}, on Schur log-concavity; and Sra's three $q$-Foulkes preprints~\cite{sra}.

This search was not a complete census of the literature: MathOverflow and Mathematics Stack Exchange were searched only up to March 2024, and Google Scholar's citation lists were not read.

{\small
\begin{thebibliography}{99}
\bibitem{GR} Á. Gutiérrez and M. H. Rosas, \emph{Partial symmetries of iterated plethysms}, Annals of Combinatorics \textbf{27} (2023), 493--518. \href{https://doi.org/10.1007/s00026-023-00652-4}{doi:10.1007/s00026-023-00652-4}. Numbering follows \href{https://arxiv.org/abs/2201.00240v2}{arXiv:2201.00240v2}.
\bibitem{thesis} Á. Gutiérrez Cáceres, \emph{Positivity conditions and hook+column sequences of plethystic coefficients}, undergraduate thesis, Universidad de Sevilla, June 2021. \href{https://idus.us.es/handle/11441/133706}{Institutional repository}.
\bibitem{positivity} Á. Gutiérrez and M. H. Rosas, \emph{Necessary conditions for the positivity of Littlewood--Richardson and plethystic coefficients}, Comptes Rendus Mathématique \textbf{361} (2023). Numbering follows \href{https://arxiv.org/abs/2109.11378v2}{arXiv:2109.11378v2}.
\bibitem{LR} T. M. Langley and J. B. Remmel, \emph{The plethysm $s_\lambda[s_\mu]$ at hook and near-hook shapes}, Electronic Journal of Combinatorics \textbf{11} (2004), R11. \href{https://doi.org/10.37236/1764}{doi:10.37236/1764}.
\bibitem{itmaia} Á. Gutiérrez, \emph{Partial symmetries of iterated plethysms of Schur functions}, slides, ITMAIA 2024, Exeter, February 2024. \href{https://drive.google.com/uc?export=download&id=1AqlTl0K5C-mCIo-A49xohpmC36v0BkzR}{Slides}.
\bibitem{lifts} Á. Gutiérrez, Á. L. Martínez, M. Szwej and M. Wildon, \emph{Plethystic lifts of $q$-binomial identities}, FPSAC 2026 extended abstract, Séminaire Lotharingien de Combinatoire \textbf{95B} (2026), Article 65. \href{https://www.mat.univie.ac.at/~slc/wpapers/FPSAC2026/65.html}{Abstract page}.
\bibitem{filtration} Á. Gutiérrez, Á. L. Martínez, M. Szwej and M. Wildon, \emph{A field-independent filtration of plethystic modules for $\mathrm{SL}_2(\mathbb{F})$ that categorifies a product rule for the Cartan subalgebra of $\mathcal{U}_q(\mathfrak{sl}_2)$}, July 7, 2026. \href{https://arxiv.org/abs/2607.06749}{arXiv:2607.06749}.
\bibitem{geometric} Á. Gutiérrez, R. Orellana, F. Saliola, A. Schilling and M. Zabrocki, \emph{A geometric and generating function approach to plethysm}, November 2025. \href{https://arxiv.org/abs/2511.02649v2}{arXiv:2511.02649v2}.
\bibitem{lee} M. Lee, \emph{Restriction coefficients for partitions with at most three columns}, Algebraic Combinatorics \textbf{9} (2026), no.~2, 543--555. \href{https://doi.org/10.5802/alco.476}{doi:10.5802/alco.476}.
\bibitem{fggk} I. Fischer, M. Gangl, Á. Gutiérrez and N. Kumari, \emph{Character expansions and affine Jacobi--Trudi identities}, October 6, 2026. \href{https://arxiv.org/abs/2610.08500v1}{arXiv:2610.08500v1}.
\bibitem{gk} Á. Gutiérrez and C. Krattenthaler, \emph{Schur log-concavity and the quantum Pascal triangle}, September 26, 2025. \href{https://arxiv.org/abs/2509.22648v1}{arXiv:2509.22648v1}.
\bibitem{sarkar} S. Sarkar, \emph{Stability of plethysm coefficients and modified polynomial induction}, October 14, 2026. \href{https://arxiv.org/abs/2610.07929v2}{arXiv:2610.07929v2}.
\bibitem{mlh} J. Ma, Y. Liu and Q. Huang, \emph{Low-twist matrix covariants of exterior powers: vanishing and modular phenomena}, August 13, 2026. \href{https://doi.org/10.5281/zenodo.21918985}{doi:10.5281/zenodo.21918985}.
\bibitem{sra} S. Sra, \emph{The $q$-Foulkes problem I: A fusion inequality for Macdonald plethysms}; \emph{II: Schur positivity with high probability}; \emph{III: On Foulkes' conjecture with a graded inner factor}, October 2026. \href{https://github.com/suvrit/preprints/tree/4b4203fe5df28499b081de5fb475bcaf82967894/preprints}{Repository}.
\end{thebibliography}
}
\end{document}
